\section{CUDA GPU并行编程：屋顶线模型分析}

\subsection{引言：性能评估的核心目标}

本章将深入讲解\textbf{屋顶线模型（Roofline Model）}。这是GPU高性能计算中至关重要的性能分析工
具，其核心目的在于：
\begin{enumerate}
    \item 科学评估应用程序的性能瓶颈；
    \item 判定应用属于\textbf{计算受限（Compute-Bound）}还是\textbf{内存受限（Memory-Bound）}；
    \item 基于分析结果选择正确的优化路径。
\end{enumerate}

在进行屋顶线分析之前，建议先掌握GPU硬件架构参数（如总线速度、显存带宽、指令吞吐量等）的基础知识。只有理解了这些底层指
标，才能准确解读分析结果。例如，在之前的向量规约（Vector Reduction）案例中，我们采用了共享内存、合并访问和
Shuffle操作等优化手段，正是因为通过屋顶线分析确认了该应用是典型的\textit{内存受限}应用。

\subsection{关键性能指标与算术强度}

\subsection{峰值性能参数}
评估性能需对比实测值与硬件理论峰值。以NVIDIA Ampere架构（A100）为例：
\begin{itemize}
    \item \textbf{显存带宽峰值（Peak Memory Bandwidth）}：约 $1.5\,\text{TB/s}$（即 $1500\,\text{GB/s}$）。
    \item \textbf{计算吞吐量峰值（Peak Compute Throughput）}：
    \begin{itemize}
        \item FP32（非Tensor Core）：$19.5\,\text{TFLOPS}$
        \item FP64：$9.7\,\text{TFLOPS}$
        \item INT32：$19.5\,\text{TOPS}$
    \end{itemize}
\end{itemize}
\textbf{注意}：仅凭单一指标（如达到1 TB/s带宽）无法判断性能优劣。高带宽可能源于过多的内存操作而缺乏有效计
算。我们需要一个综合指标来平衡二者。

\subsection{算术强度（Arithmetic Intensity, AI）}
算术强度是屋顶线模型的核心度量，定义为每字节内存访问所执行的浮点运算次数：
\begin{equation}
    \text{AI} = \frac{\text{指令吞吐量 (FLOPS/s)}}
{\text{内存带宽(Bytes/s)}} \quad [\text{单位: FLOP/Byte}]
\end{equation}
\begin{itemize}
    \item \textbf{高AI}：计算量远大于数据搬运量 $\rightarrow$ 倾向于计算受限。
    \item \textbf{低AI}：数据搬运量远大于计算量 $\rightarrow$ 倾向于内存受限。
\end{itemize}

\subsection{可达性能公式}
应用程序的可达性能（Attainable Performance）受限于硬件峰值和当前算术强度下的带宽上限，取二者最小值：
\begin{equation}
    \text{Performance} = \min \left( \text{Peak Compute}
,\ \quad \text{AI} \times \text{Peak Bandwidth} \right)
\end{equation}
该公式构成了屋顶线图的数学基础：当 $\text{AI} \times \text{Bandwidth} < \text{Peak Compute}$ 时，
性能随AI线性增长（内存墙）；反之则达到平台期（计算墙）。

\subsection{屋顶线图解读与转折点}

\subsection{坐标轴定义}
\begin{itemize}
    \item \textbf{X轴（对数坐标）}：算术强度（AI），单位 FLOP/Byte。
    \item \textbf{Y轴（对数坐标）}：性能，单位 GFLOPS 或 TFLOPS。
\end{itemize}

\subsection{转折点（Ridge Point）}
转折点是内存受限区与计算受限区的分界点，对应硬件同时达到带宽峰值和计算峰值的状态：
\begin{equation}
    \text{AI}_{\text{ridge}} = \frac
{\text{Peak Compute(FLOPS/s)}}
{\text{Peak Bandwidth (Bytes/s)}}
\end{equation}
以A100 FP32为例：$\text{AI}_{\text{ridge}} = 19500 / 1500 =1
3\,\text{FLOP/Byte}$。
\begin{itemize}
    \item \textbf{左侧区域}（$\text{AI} < \text{AI}_{\text{ridge}}$）：\textbf{内存受限区}。优化重点：合并访问、共享内存、L1/L2缓存利用率、预取。
    \item \textbf{右侧区域}（$\text{AI} > \text{AI}_{\text{ridge}}$）：\textbf{计算受限区}。优化重点：指令级并行、消除依赖、Fast Math、Tensor Core、寄存器优化。
\end{itemize}

\subsection{多级存储屋顶线}
实际分析中需区分不同存储层级，因其带宽差异巨大：
\begin{itemize}
    \item \textbf{DRAM Roofline}：使用全局显存峰值带宽（最常用基准）。
    \item \textbf{L2 Cache Roofline}：使用L2缓存峰值带宽（转折点右移）。
    \item \textbf{L1/Shared Mem Roofline}：使用L1/SMEM峰值带宽（转折点进一步右移）。
\end{itemize}
同一应用在DRAM层面可能是内存受限，但在L1层面可能已进入计算受限区。优化目标是将数据点向右上方推移，逼近或超越转折点
。

\subsection{实战案例：向量规约的屋顶线分析}

\subsection{基准版本分析}
使用Nsight Compute对向量规约基准版本（01\_reduce）进行分析：
\begin{lstlisting}
nvcc -O3 -o reduce_01 reduce_01.cu
ncu --launch-mode profile ./reduce_01
\end{lstlisting}
在Nsight Compute的``Roofline''面板中选择``FP32(Non-Tensor)''视图。对于RTX 3060等消费级卡，转折点AI约为30FLOP/Byte。基准版本的实测点位于转折点\textbf{左侧远处}
（AI$\approx$0.12），确认为严重内存受限。
\subsection{优化版本对比} 对优化版本（03\_reduce\_shared\_unroll）进行相同分析：
\begin{itemize}
    \item \textbf{基准版}：AI = 0.12 FLOP/B, 性能 $\approx$ 15 GFLOPS
    \item \textbf{优化版}：AI = 0.66 FLOP/B, 性能 $\approx$ 177 GFLOPS
\end{itemize}
\textbf{分析结论}：
\begin{enumerate}
    \item 算术强度提升约 \textbf{5.5倍}（$0.66 / 0.12$）；
    \item 绝对性能提升约 \textbf{11.8倍}（$177 / 15$）；
    \item 数据点在屋顶线图中显著向右上方移动，但仍处于转折点左侧，说明仍有内存优化空间（如后续Shuffle优化可进一步提升AI）。
\end{enumerate}

\subsection{总结与优化策略指引}

屋顶线模型为GPU性能优化提供了量化导航：
\begin{table}[h]
\centering
\caption{基于屋顶线位置的优化策略}
\begin{tabular}{@{}lll@{}}
\toprule
\textbf{区域} & \textbf{特征} & \textbf{推荐优化手段} \\ \midrule
内存受限区 & AI < Ridge Point & 合并访存、Shared Memory、向量化加载、\\
& & L1/L2 Cache复用、减少冗余读写 \\ \midrule
计算受限区 & AI > Ridge Point & ILP/TLP提升、消除分支发散、\\
& & Fast Math Intrinsics、Tensor Core、\\
& & 算法复杂度优化 \\ \midrule
转折点附近 & AI $\approx$ Ridge Point & 硬件已充分利用，考虑算法级重构 \\
\bottomrule
\end{tabular}
\end{table}

\textbf{核心原则}：优化的终极目标是使应用的数据点在屋顶线图中持续向右上方迁移，直至逼近硬件物理极限构成的``屋
顶''。每次优化后都应重新运行Nsight Compute验证数据点位移，形成``测量$\to$分析$\to$优化$\to$验证''的闭环迭代。
